\setcounter{chapter}{3}
\chapter*{Activités Logiques}
\section*{Traduction du Langage Naturel en Logique Propositionnelle}
\subsection*{Exercice 1 : Logique propositionnelle}


Traduisez l'affirmation suivante en utilisant la notation de la logique propositionnelle :
\begin{center}
"Si il neige (N) et que le vent souffle (V), alors il fera froid (F) ; cependant, même s'il fait froid, on peut ne pas avoir ni neige, ni vent."
\end{center}
Connecteurs Logiques Clés :
\begin{itemize}
    \item ET (conjonction) : \(\land\)
    \item OU (disjonction) : \(\lor\)
    \item NON (négation) : \(\neg\)
    \item IMPLIQUE (implication) : \(\rightarrow\) (ou \(\Rightarrow\))
\end{itemize}
Variables :
\begin{itemize}
    \item N
    \item V
    \item F
\end{itemize}

\subsection*{Exercice 2 : La Clause Contractuelle}

Traduisez l'affirmation suivante en utilisant la notation de la logique propositionnelle.
\begin{center}
    "Vous obtiendrez une réduction (R) à la seule condition que vous achetiez le forfait premium (P), mais si vous achetez le forfait premium, vous devez également fournir un dépôt de garantie (D)."
\end{center}
Variables :
\begin{itemize}
    \item Vous obtiendrez une réduction.
    \item Vous achetez le forfait premium.
    \item Vous fournissez un dépôt de garantie.
\end{itemize}

Connecteurs Logiques Clés\\
\begin{itemize}
    \item "à la seule condition que" : Il s'agit d'une condition nécessaire, pas d'une équivalence : "A à la seule condition que B" se traduit par \( A \rightarrow B \).
\end{itemize}
\begin{Comprendre}
  L'expression "si et seulement si" (\( \Leftrightarrow \)) est la plus stricte. "A à la seule condition que B" signifie que \( A \) ne peut être vrai que si \( B \) est vrai, soit \( A \rightarrow B \).  
\end{Comprendre}

\subsection*{Exercice 3 : La Condition Météo}
Traduisez l'affirmation suivante en logique propositionnelle.
\begin{center}
"Le pique-nique sera annulé (A) s'il pleut (P) ou s'il y a du vent fort (V). De plus, s'il ne pleut pas et qu'il n'y a pas de vent fort, le pique-nique aura lieu, à condition que la température soit clémente (T)."
\end{center}
Variables :
\begin{itemize}
    \item Le pique-nique sera annulé :
    \item Il pleut :
    \item Il y a du vent fort :
    \item La température est clémente :
\end{itemize}


Connecteurs Logiques Clés
\begin{itemize}
    \item  "s'il..." : Souvent une implication ($\rightarrow$).
    \item "ou" : Disjonction ($\lor$).
    \item  "ne... pas" : Négation ($\neg$).
    \item "à condition que" : Implication ($\rightarrow$).
\end{itemize}
    


\section*{Réflexion en groupes }
\subsection*{Activité 1 : Les Cartes Logiques (Logique Propositionnelle)}

\textbf{Le Principe}\\
Nous allons résoudre l’énigme des Quatre Cartes de Wason, un célèbre test de raisonnement logique.\\
Vous avez quatre cartes. Chaque carte a une lettre d'un côté et un nombre de l'autre.\\

%
\textbf{L'Énoncé}\\
S'il y a un A sur une face, alors il y a un 4 sur l'autre face.
\textbf{Quelles sont les seules cartes qu'il est absolument nécessaire de retourner pour vérifier si la règle est respectée ou non ?}\\
Les Cartes :
\begin{enumerate}
    \item A (Lettre)
    \item K (Lettre)
    \item 4 (Nombre)
    \item 7 (Nombre)
\end{enumerate}

Aide : Réfléchir à ce qu’on peut prouver en retournant chacune des cartes.
\begin{itemize}
    \item Carte A 
    \item Carte K 
    \item Carte 4 
    \item Carte 7 
\end{itemize}


%
\subsection*{Activité 2 : Logique des Instructions (Algorithmique et Binaire)}
Vous devez programmer un robot (imaginaire) pour atteindre un objectif précis en utilisant un jeu d'instructions très limité.\\

%
\textbf{Le Défi - Simulateur de Code}\\
Vous avez une grille avec un point de départ où se trouve le Robot et un point d'arrivée.\\
Instructions Autorisées :
\begin{itemize}
    \item AVANCER (se déplace d'une case dans la direction actuelle)
    \item TOURNER\_DROITE
    \item SI\_BLOCAGE (instruction conditionnelle)
    \item FIN\_SI
    \item RÉPÉTER N (boucle)
    \item FIN\_RÉPÉTER
\end{itemize}

%
\textbf{L’Enoncé}\\
 Écrire le programme le plus court pour faire passer le robot du coin supérieur gauche au coin inférieur droit, en évitant un obstacle au centre.

